% ============================================================================
% CAPÍTULO 12: PHYSIOLOGICAL SYSTEMS - HEART
% Bioinformática para Biologia de Sistemas
% ============================================================================

% ============================================================================
% TEMPLATE BASE PARA APRESENTAÇÕES EM BEAMER
% Bioinformática para Biologia de Sistemas
% ============================================================================

\documentclass[12pt, aspectratio=169]{beamer}

% ============================================================================
% PACOTES ESSENCIAIS
% ============================================================================
\usepackage[utf8]{inputenc}
\usepackage[T1]{fontenc}
\usepackage[brazilian]{babel}
\usepackage{amsmath, amssymb, amsthm}
\usepackage{graphicx}
\usepackage{xcolor}
\usepackage{tikz}
\usetikzlibrary{arrows.meta, shapes, positioning, calc, decorations.pathreplacing, decorations.shapes, decorations.pathmorphing}
\usepackage{listings}
\usepackage{booktabs}
\usepackage{multirow}
\usepackage{hyperref}
\usepackage{chemfig}
\usepackage{chemarr}

% ============================================================================
% CONFIGURAÇÃO DO TEMA
% ============================================================================
\usetheme{Madrid}
\usecolortheme{default}

% Cores personalizadas
\definecolor{azulescuro}{RGB}{0, 51, 102}
\definecolor{azulclaro}{RGB}{51, 153, 255}
\definecolor{verdeacido}{RGB}{102, 204, 0}
\definecolor{laranjacel}{RGB}{255, 153, 51}
\definecolor{roxodna}{RGB}{153, 51, 255}

\setbeamercolor{structure}{fg=azulescuro}
\setbeamercolor{palette primary}{bg=azulescuro,fg=white}
\setbeamercolor{palette secondary}{bg=azulclaro,fg=white}
\setbeamercolor{palette tertiary}{bg=azulescuro,fg=white}
\setbeamercolor{palette quaternary}{bg=azulclaro,fg=white}
\setbeamercolor{block title}{bg=azulescuro,fg=white}
\setbeamercolor{block body}{bg=azulclaro!10,fg=black}
\setbeamercolor{block title alerted}{bg=red!80,fg=white}
\setbeamercolor{block body alerted}{bg=red!10,fg=black}
\setbeamercolor{block title example}{bg=verdeacido!80,fg=white}
\setbeamercolor{block body example}{bg=verdeacido!10,fg=black}

% ============================================================================
% CONFIGURAÇÃO DE CÓDIGO
% ============================================================================
\lstset{
    language=Python,
    basicstyle=\ttfamily\small,
    keywordstyle=\color{azulescuro}\bfseries,
    commentstyle=\color{verdeacido},
    stringstyle=\color{laranjacel},
    numbers=left,
    numberstyle=\tiny\color{gray},
    stepnumber=1,
    numbersep=5pt,
    backgroundcolor=\color{white},
    showspaces=false,
    showstringspaces=false,
    showtabs=false,
    frame=single,
    rulecolor=\color{black},
    tabsize=2,
    captionpos=b,
    breaklines=true,
    breakatwhitespace=false,
    escapeinside={\%*}{*)},
    xleftmargin=10pt,
    xrightmargin=10pt
}

% Definir estilo para R
\lstdefinestyle{Rstyle}{
    language=R,
    basicstyle=\ttfamily\small,
    keywordstyle=\color{azulescuro}\bfseries,
    commentstyle=\color{verdeacido},
    stringstyle=\color{laranjacel}
}

% Definir estilo para MATLAB
\lstdefinestyle{MATLABstyle}{
    language=Matlab,
    basicstyle=\ttfamily\small,
    keywordstyle=\color{azulescuro}\bfseries,
    commentstyle=\color{verdeacido},
    stringstyle=\color{laranjacel}
}

% ============================================================================
% COMANDOS PERSONALIZADOS
% ============================================================================

% Comando para equações destacadas
\newcommand{\destaque}[1]{%
    \begin{alertblock}{}
        \begin{center}
            #1
        \end{center}
    \end{alertblock}
}

% Comando para definições
\newcommand{\definicao}[2]{%
    \begin{block}{Definição: #1}
        #2
    \end{block}
}

% Comando para exemplos
\newcommand{\exemplo}[2]{%
    \begin{exampleblock}{Exemplo: #1}
        #2
    \end{exampleblock}
}

% Comando para observações importantes
\newcommand{\importante}[1]{%
    \begin{alertblock}{Importante!}
        #1
    \end{alertblock}
}

% Comandos matemáticos úteis
\newcommand{\R}{\mathbb{R}}
\newcommand{\N}{\mathbb{N}}
\newcommand{\Z}{\mathbb{Z}}
\newcommand{\dif}{\mathrm{d}}
\newcommand{\parcial}[2]{\frac{\partial #1}{\partial #2}}
\newcommand{\derivada}[2]{\frac{\dif #1}{\dif #2}}
\newcommand{\vect}[1]{\boldsymbol{#1}}

% Comandos biológicos
\newcommand{\gene}[1]{\textit{#1}}
\newcommand{\proteina}[1]{\textsc{#1}}
\newcommand{\especie}[1]{\textit{#1}}

% ============================================================================
% CONFIGURAÇÃO DE RODAPÉ
% ============================================================================
\setbeamertemplate{footline}{
  \leavevmode%
  \hbox{%
  \begin{beamercolorbox}[wd=.33\paperwidth,ht=2.25ex,dp=1ex,center]{author in head/foot}%
    \usebeamerfont{author in head/foot}\insertshortauthor
  \end{beamercolorbox}%
  \begin{beamercolorbox}[wd=.34\paperwidth,ht=2.25ex,dp=1ex,center]{title in head/foot}%
    \usebeamerfont{title in head/foot}\insertshorttitle
  \end{beamercolorbox}%
  \begin{beamercolorbox}[wd=.33\paperwidth,ht=2.25ex,dp=1ex,right]{date in head/foot}%
    \usebeamerfont{date in head/foot}\insertshortdate{}\hspace*{2em}
    \insertframenumber{} / \inserttotalframenumber\hspace*{2ex}
  \end{beamercolorbox}}%
  \vskip0pt%
}

% ============================================================================
% CONFIGURAÇÃO DE NAVEGAÇÃO
% ============================================================================
\setbeamertemplate{navigation symbols}{}

% ============================================================================
% CONFIGURAÇÃO DE ITEMIZE/ENUMERATE
% ============================================================================
\setbeamertemplate{itemize items}[circle]
\setbeamertemplate{enumerate items}[default]

% ============================================================================
% FIM DO TEMPLATE
% ============================================================================


% Pacotes adicionais
\usepackage{chemfig}

% ============================================================================
% INFORMAÇÕES DO DOCUMENTO
% ============================================================================
\title[Cap. 12: Heart Physiology]{Capítulo 12: Physiological Systems - Heart}
\subtitle{Bioinformática para Biologia de Sistemas}
\author{Seu Nome}
\institute{Sua Instituição}
\date{\today}

\begin{document}

% ============================================================================
% SLIDE DE TÍTULO
% ============================================================================
\begin{frame}
\titlepage
\end{frame}

% ============================================================================
% SUMÁRIO
% ============================================================================
\begin{frame}{Sumário}
\tableofcontents
\end{frame}

% ============================================================================
% SEÇÃO 1: INTRODUÇÃO
% ============================================================================
\section{Introdução}

\begin{frame}{O Coração como Sistema Complexo}
\definicao{Sistema Cardíaco}{
O coração é um sistema fisiológico complexo que integra eletrofisiologia, mecânica e regulação para bombear sangue continuamente através do organismo.
}

\vspace{0.3cm}

\begin{columns}
\column{0.5\textwidth}
\textbf{Características:}
\begin{itemize}
    \item Bomba muscular rítmica
    \item Automaticidade intrínseca
    \item Acoplamento elétrico-mecânico
    \item Regulação autônoma
    \item Adaptação à demanda metabólica
\end{itemize}

\column{0.5\textwidth}
\begin{center}
\begin{tikzpicture}[scale=0.8]
    % Heart outline simplified
    \draw[thick, fill=red!30] (0,0) .. controls (0,1) and (1,1.5) .. (1.5,1.5)
        .. controls (2,1.5) and (3,1) .. (3,0)
        .. controls (2.5,-1) and (2,-2) .. (1.5,-2.5)
        .. controls (1,-2) and (0.5,-1) .. (0,0);

    % Heart chambers indication
    \draw[thick] (1.5,1) -- (1.5,-1.5);
    \draw[thick] (0.7,0.5) -- (2.3,0.5);

    % Labels
    \node[font=\tiny] at (0.7,1) {AE};
    \node[font=\tiny] at (2.3,1) {AD};
    \node[font=\tiny] at (0.7,-0.5) {VE};
    \node[font=\tiny] at (2.3,-0.5) {VD};
\end{tikzpicture}
\end{center}
\textbf{Frequência:} 60-100 bpm\\
\textbf{Débito:} 5 L/min (repouso)
\end{columns}
\end{frame}

\begin{frame}{Níveis de Organização Cardíaca}
\begin{center}
\begin{tikzpicture}[node distance=2cm, auto,
    box/.style={rectangle, draw, text width=3cm, text centered, rounded corners, minimum height=0.8cm, font=\small}]

    \node[box, fill=roxodna!30] (canal) {Canais Iônicos};
    \node[box, fill=laranjacel!30, below of=canal] (celula) {Cardiomiócito};
    \node[box, fill=verdeacido!30, below of=celula] (tecido) {Tecido Cardíaco};
    \node[box, fill=azulclaro!30, below of=tecido] (orgao) {Órgão (Coração)};
    \node[box, fill=yellow!30, below of=orgao] (sistema) {Sistema Cardiovascular};

    \draw[->, thick] (canal) -- (celula);
    \draw[->, thick] (celula) -- (tecido);
    \draw[->, thick] (tecido) -- (orgao);
    \draw[->, thick] (orgao) -- (sistema);

    % Annotations
    \node[right, font=\tiny, text width=3cm] at (4,0) {Na$^+$, K$^+$, Ca$^{2+}$};
    \node[right, font=\tiny, text width=3cm] at (4,-2) {Potencial de ação};
    \node[right, font=\tiny, text width=3cm] at (4,-4) {Propagação elétrica};
    \node[right, font=\tiny, text width=3cm] at (4,-6) {Contração coordenada};
    \node[right, font=\tiny, text width=3cm] at (4,-8) {Perfusão tecidual};
\end{tikzpicture}
\end{center}
\end{frame}

\begin{frame}{Visão Geral do Capítulo}
\begin{block}{Objetivos de Aprendizagem}
\begin{itemize}
    \item Compreender a anatomia e fisiologia cardíaca
    \item Dominar a eletrofisiologia do cardiomiócito
    \item Estudar modelos matemáticos de células cardíacas
    \item Entender o acoplamento excitação-contração
    \item Explorar mecanismos de arritmias
    \item Aplicar modelagem computacional ao coração
\end{itemize}
\end{block}

\vspace{0.3cm}

\importante{
O coração é um dos sistemas fisiológicos mais bem modelados matematicamente, com implicações clínicas diretas para diagnóstico e terapia de doenças cardiovasculares.
}
\end{frame}

% ============================================================================
% SEÇÃO 2: ANATOMIA E FISIOLOGIA CARDÍACA
% ============================================================================
\section{Anatomia e Fisiologia Cardíaca}

\begin{frame}{Estrutura Anatômica do Coração}
\begin{columns}
\column{0.5\textwidth}
\textbf{Câmaras:}
\begin{itemize}
    \item \textbf{Átrios:} direito e esquerdo\\
    (recepção de sangue)
    \item \textbf{Ventrículos:} direito e esquerdo\\
    (bombeamento)
\end{itemize}

\vspace{0.3cm}

\textbf{Valvas:}
\begin{itemize}
    \item \textbf{Atrioventriculares:}\\
    Tricúspide (AD-VD)\\
    Mitral (AE-VE)
    \item \textbf{Semilunares:}\\
    Pulmonar (VD)\\
    Aórtica (VE)
\end{itemize}

\column{0.5\textwidth}
\begin{center}
\begin{tikzpicture}[scale=0.7]
    % Heart chambers
    % Right atrium
    \draw[thick, fill=azulclaro!20] (0,1) rectangle (1.5,3);
    \node at (0.75,2) {\tiny AD};

    % Right ventricle
    \draw[thick, fill=azulclaro!40] (0,-1.5) -- (0,0.5) -- (1.5,0.5) -- (1.8,-1.5) -- cycle;
    \node at (0.75,-0.5) {\tiny VD};

    % Left atrium
    \draw[thick, fill=red!20] (2,1) rectangle (3.5,3);
    \node at (2.75,2) {\tiny AE};

    % Left ventricle
    \draw[thick, fill=red!40] (1.7,-1.5) -- (2,0.5) -- (3.5,0.5) -- (3.5,-1.5) -- cycle;
    \node at (2.75,-0.5) {\tiny VE};

    % Valves
    \draw[thick, red] (0.5,0.5) -- (1,0.5);
    \node[left, font=\tiny] at (0.3,0.5) {Tric};

    \draw[thick, red] (2.5,0.5) -- (3,0.5);
    \node[right, font=\tiny] at (3.3,0.5) {Mitral};

    % Great vessels
    \draw[->, ultra thick, azulclaro] (0.75,-1.5) -- (0.75,-2.5) node[below, font=\tiny] {Pulmonar};
    \draw[->, ultra thick, red] (2.75,-1.5) -- (2.75,-2.5) node[below, font=\tiny] {Aorta};

    % Septum
    \draw[thick, dashed] (1.75,3) -- (1.75,-1.5);
    \node[right, font=\tiny] at (1.75,1.5) {Septo};
\end{tikzpicture}
\end{center}
\end{columns}

\vspace{0.3cm}

\begin{exampleblock}{Circulação}
\textbf{Direito:} sangue venoso $\rightarrow$ pulmões (oxigenação)\\
\textbf{Esquerdo:} sangue arterial $\rightarrow$ corpo (perfusão sistêmica)
\end{exampleblock}
\end{frame}

\begin{frame}{Ciclo Cardíaco}
\begin{block}{Fases do Ciclo}
\begin{enumerate}
    \item \textbf{Diástole:} relaxamento e enchimento
    \begin{itemize}
        \item Enchimento ventricular passivo
        \item Contração atrial (sístole atrial)
    \end{itemize}
    \item \textbf{Sístole:} contração e ejeção
    \begin{itemize}
        \item Contração isovolumétrica
        \item Ejeção ventricular
        \item Relaxamento isovolumétrico
    \end{itemize}
\end{enumerate}
\end{block}

\begin{center}
\begin{tikzpicture}[scale=0.8]
    % Pressure-volume loop
    \draw[->] (0,0) -- (5,0) node[right] {Volume (mL)};
    \draw[->] (0,0) -- (0,4) node[above] {Pressão (mmHg)};

    % Loop
    \draw[thick, azulescuro, ->] (1,0.5) -- node[below] {\tiny Enchimento} (3.5,0.5);
    \draw[thick, azulescuro, ->] (3.5,0.5) -- node[right] {\tiny Contração isovolumétrica} (3.5,3);
    \draw[thick, azulescuro, ->] (3.5,3) -- node[above] {\tiny Ejeção} (1,3);
    \draw[thick, azulescuro, ->] (1,3) -- node[left] {\tiny Relaxamento isovolumétrico} (1,0.5);

    % Labels
    \node[below] at (1,0) {\tiny ESV};
    \node[below] at (3.5,0) {\tiny EDV};
\end{tikzpicture}
\end{center}

\importante{
Duração típica: $\sim$0.8 s a 75 bpm (diástole $\sim$0.5 s, sístole $\sim$0.3 s)
}
\end{frame>

\begin{frame}{Sistema de Condução Elétrica}
\begin{columns}
\column{0.5\textwidth}
\textbf{Componentes:}
\begin{enumerate}
    \item \textbf{Nó Sinoatrial (SA):}\\
    Marca-passo natural\\
    60-100 bpm

    \item \textbf{Nó Atrioventricular (AV):}\\
    Atraso de condução\\
    $\sim$0.1 s

    \item \textbf{Feixe de His:}\\
    Condução rápida

    \item \textbf{Ramos (direito/esquerdo):}\\
    Distribuição ventricular

    \item \textbf{Fibras de Purkinje:}\\
    Ativação endocárdio
\end{enumerate}

\column{0.5\textwidth}
\begin{center}
\begin{tikzpicture}[scale=0.8]
    % Heart outline
    \draw[thick] (0,0) ellipse (1.5cm and 2cm);

    % SA node
    \node[circle, draw, fill=red!50, minimum size=0.4cm] (sa) at (0.5,1.5) {};
    \node[right, font=\tiny] at (0.7,1.5) {Nó SA};

    % Atrial conduction
    \draw[->, thick, red] (sa) -- (-0.5,1);

    % AV node
    \node[circle, draw, fill=yellow!50, minimum size=0.4cm] (av) at (0,0.3) {};
    \node[right, font=\tiny] at (0.2,0.3) {Nó AV};

    % Bundle of His
    \draw[->, thick, verdeacido] (av) -- (0,-0.5);
    \node[right, font=\tiny] at (0.2,-0.5) {Feixe His};

    % Bundle branches
    \draw[->, thick, azulclaro] (0,-0.5) -- (-0.8,-1.5);
    \draw[->, thick, azulclaro] (0,-0.5) -- (0.8,-1.5);

    % Purkinje fibers
    \draw[thick, roxodna, decoration={coil, amplitude=1pt}, decorate]
        (-0.8,-1.5) -- (-1.2,-1.8);
    \draw[thick, roxodna, decoration={coil, amplitude=1pt}, decorate]
        (0.8,-1.5) -- (1.2,-1.8);
    \node[below, font=\tiny] at (0,-2) {Fibras Purkinje};
\end{tikzpicture}
\end{center}
\end{columns}

\vspace{0.3cm}

\destaque{
Velocidade de condução: Purkinje (4 m/s) $>$ His (1-1.5 m/s) $>$ AV (0.05 m/s)
}
\end{frame>

\begin{frame}{Eletrocardiograma (ECG)}
\begin{columns}
\column{0.5\textwidth}
\textbf{Ondas do ECG:}
\begin{itemize}
    \item \textbf{Onda P:} despolarização atrial
    \item \textbf{Complexo QRS:} despolarização ventricular
    \item \textbf{Onda T:} repolarização ventricular
    \item \textbf{Intervalo PR:} 0.12-0.20 s
    \item \textbf{Intervalo QT:} 0.35-0.44 s
    \item \textbf{Segmento ST:} platô ventricular
\end{itemize}

\vspace{0.3cm}

\begin{alertblock}{Significado Clínico}
\begin{itemize}
    \item Arritmias
    \item Isquemia/infarto
    \item Distúrbios de condução
    \item Efeitos de drogas
\end{itemize}
\end{alertblock}

\column{0.5\textwidth}
\begin{center}
\begin{tikzpicture}[scale=0.8]
    % ECG trace
    \draw[->] (0,0) -- (8,0) node[right] {Tempo (s)};
    \draw[->] (0,-1) -- (0,3) node[above] {mV};

    % Baseline
    \draw[help lines] (0,0) -- (8,0);

    % P wave
    \draw[thick, azulescuro] (1,0) .. controls (1.2,0.3) and (1.4,0.3) .. (1.6,0);
    \node[above, font=\tiny] at (1.3,0.4) {P};

    % QRS complex
    \draw[thick, azulescuro] (2.5,0) -- (2.7,-0.5) -- (3,2.5) -- (3.3,-0.7) -- (3.5,0);
    \node[below, font=\tiny] at (2.7,-0.6) {Q};
    \node[above, font=\tiny] at (3,2.7) {R};
    \node[below, font=\tiny] at (3.3,-0.8) {S};

    % ST segment
    \draw[thick, azulescuro] (3.5,0) -- (4.5,0);

    % T wave
    \draw[thick, azulescuro] (4.5,0) .. controls (5,0.8) and (5.5,0.8) .. (6,0);
    \node[above, font=\tiny] at (5.2,0.9) {T};

    % Intervals
    \draw[<->, red] (1.3,-1.5) -- node[below, font=\tiny] {Intervalo PR} (2.5,-1.5);
    \draw[<->, verdeacido] (2.5,-2) -- node[below, font=\tiny] {Intervalo QT} (6,-2);
\end{tikzpicture}
\end{center>
\end{columns}
\end{frame}

\begin{frame}{Circulação Coronariana}
\begin{block}{Artérias Coronárias}
\begin{itemize}
    \item \textbf{Coronária Direita (CD):} VD, parede inferior VE, nó SA/AV
    \item \textbf{Artéria Descendente Anterior Esquerda (DAE):} septo, parede anterior VE
    \item \textbf{Artéria Circunflexa (Cx):} parede lateral e posterior VE
\end{itemize}
\end{block}

\vspace{0.3cm}

\begin{columns}
\column{0.5\textwidth}
\textbf{Características:}
\begin{itemize}
    \item Fluxo principalmente diastólico
    \item Autorregulação metabólica
    \item Alta extração de O$_2$ ($\sim$75\%)
    \item Reserva vasodilatadora limitada
\end{itemize}

\column{0.5\textwidth}
\begin{exampleblock}{Fluxo Coronariano}
\begin{equation*}
Q = \frac{\Delta P}{R_{cor}}
\end{equation*}
onde:
\begin{itemize}
    \item $Q$ = fluxo ($\sim$250 mL/min)
    \item $\Delta P$ = gradiente de pressão
    \item $R_{cor}$ = resistência coronariana
\end{itemize}
\end{exampleblock}
\end{columns}

\vspace{0.3cm}

\importante{
Isquemia ocorre quando demanda de O$_2$ excede oferta (estenose, espasmo, aumento de demanda)
}
\end{frame}

% ============================================================================
% SEÇÃO 3: ELETROFISIOLOGIA DO CARDIOMIÓCITO
% ============================================================================
\section{Eletrofisiologia do Cardiomiócito}

\begin{frame}{Potencial de Ação Cardíaco}
\begin{columns}
\column{0.5\textwidth}
\textbf{Fases (Ventricular):}
\begin{enumerate}
    \item[\textbf{0:}] Despolarização rápida\\
    Entrada de Na$^+$ (I$_{Na}$)

    \item[\textbf{1:}] Repolarização inicial\\
    Saída de K$^+$ (I$_{to}$)

    \item[\textbf{2:}] Platô\\
    Entrada Ca$^{2+}$ (I$_{Ca,L}$)\\
    Saída K$^+$ (I$_{Kr}$, I$_{Ks}$)

    \item[\textbf{3:}] Repolarização final\\
    Saída de K$^+$ (I$_{K1}$)

    \item[\textbf{4:}] Repouso\\
    Potencial de membrana estável
\end{enumerate}

\column{0.5\textwidth}
\begin{center}
\begin{tikzpicture}[scale=0.8]
    % Axes
    \draw[->] (0,0) -- (6,0) node[right, font=\tiny] {Tempo (ms)};
    \draw[->] (0,-1) -- (0,4) node[above, font=\tiny] {V$_m$ (mV)};

    % Resting potential
    \draw[thick, azulescuro] (0,0) -- (1,0);
    \node[left, font=\tiny] at (0,0) {-90};

    % Phase 0 - upstroke
    \draw[thick, azulescuro] (1,0) -- (1.3,3.5);
    \node[right, font=\tiny, red] at (1.3,1.5) {0};
    \node[left, font=\tiny] at (1.3,3.5) {+30};

    % Phase 1 - early repol
    \draw[thick, azulescuro] (1.3,3.5) -- (1.7,2.8);
    \node[right, font=\tiny, red] at (1.5,3) {1};

    % Phase 2 - plateau
    \draw[thick, azulescuro] (1.7,2.8) -- (3.5,2.5);
    \node[above, font=\tiny, red] at (2.5,2.7) {2};

    % Phase 3 - repolarization
    \draw[thick, azulescuro] (3.5,2.5) -- (4.5,0);
    \node[right, font=\tiny, red] at (4,1.2) {3};

    % Phase 4 - resting
    \draw[thick, azulescuro] (4.5,0) -- (6,0);
    \node[below, font=\tiny, red] at (5.2,-0.2) {4};

    % Ion currents
    \draw[->, verdeacido, very thick] (1.1,0.5) -- (1.1,1.5) node[left, font=\tiny] {Na$^+$};
    \draw[->, laranjacel, very thick] (2.5,2.2) -- (2.5,1.7) node[below, font=\tiny] {Ca$^{2+}$};
    \draw[->, roxodna, very thick] (4,1.5) -- (4,0.5) node[right, font=\tiny] {K$^+$};
\end{tikzpicture}
\end{center}

\textbf{Duração:} 200-400 ms\\
\textbf{Amplitude:} $\sim$120 mV
\end{columns}

\vspace{0.3cm}

\importante{
O platô prolongado (Fase 2) previne tetania cardíaca e garante enchimento ventricular adequado
}
\end{frame}

\begin{frame}{Canais Iônicos Cardíacos}
\begin{block}{Principais Correntes Iônicas}
\begin{center}
\begin{tabular}{llll}
\toprule
\textbf{Corrente} & \textbf{Canal} & \textbf{Fase} & \textbf{Função} \\
\midrule
I$_{Na}$ & Nav1.5 & 0 & Despolarização \\
I$_{to}$ & Kv4.3 & 1 & Repol. inicial \\
I$_{Ca,L}$ & Cav1.2 & 2 & Platô, entrada Ca$^{2+}$ \\
I$_{Kr}$ & hERG & 2-3 & Repol. (rápido) \\
I$_{Ks}$ & KCNQ1/KCNE1 & 2-3 & Repol. (lento) \\
I$_{K1}$ & Kir2.x & 3-4 & Repol. final, repouso \\
I$_{f}$ & HCN & 4 & Automaticidade (SA) \\
\bottomrule
\end{tabular}
\end{center}
\end{block}

\vspace{0.3cm}

\begin{columns}
\column{0.5\textwidth}
\begin{exampleblock}{Gating}
Canais possuem gates (comportas):
\begin{itemize}
    \item \textbf{m:} ativação
    \item \textbf{h:} inativação rápida
    \item \textbf{j:} inativação lenta
\end{itemize}
\end{exampleblock}

\column{0.5\textwidth}
\begin{alertblock}{Canalopatias}
Mutações em canais causam:
\begin{itemize}
    \item Síndrome do QT longo
    \item Síndrome de Brugada
    \item Taquicardia ventricular
    \item Morte súbita
\end{itemize}
\end{alertblock}
\end{columns}
\end{frame}

\begin{frame}{Potencial de Membrana em Repouso}
\begin{block}{Equação de Goldman-Hodgkin-Katz}
Para múltiplos íons:
\begin{equation*}
V_m = \frac{RT}{F}\ln\left(\frac{P_K[K^+]_o + P_{Na}[Na^+]_o + P_{Cl}[Cl^-]_i}{P_K[K^+]_i + P_{Na}[Na^+]_i + P_{Cl}[Cl^-]_o}\right)
\end{equation*}
\end{block}

\vspace{0.3cm}

\begin{columns}
\column{0.5\textwidth}
\textbf{Concentrações Típicas (mM):}
\begin{center}
\begin{tabular}{lcc}
\toprule
\textbf{Íon} & \textbf{Intra} & \textbf{Extra} \\
\midrule
K$^+$ & 140 & 5 \\
Na$^+$ & 10 & 145 \\
Ca$^{2+}$ & 0.0001 & 2 \\
Cl$^-$ & 10 & 110 \\
\bottomrule
\end{tabular}
\end{center}

\column{0.5\textwidth}
\textbf{Potenciais de Equilíbrio:}
\begin{align*}
E_K &\approx -90 \text{ mV} \\
E_{Na} &\approx +70 \text{ mV} \\
E_{Ca} &\approx +130 \text{ mV} \\
E_{Cl} &\approx -70 \text{ mV}
\end{align*}

Em repouso (Fase 4):\\
$V_m \approx -85$ a $-90$ mV\\
(próximo de $E_K$)
\end{columns}

\vspace{0.3cm}

\destaque{
$V_m$ em repouso dominado por permeabilidade ao K$^+$ via I$_{K1}$
}
\end{frame>

\begin{frame}{Mecanismo do Platô (Fase 2)}
\begin{columns}
\column{0.5\textwidth}
\textbf{Balanço de Correntes:}

\vspace{0.3cm}

\textbf{Entrada de carga (+):}
\begin{itemize}
    \item I$_{Ca,L}$: corrente de Ca$^{2+}$ tipo-L
    \item Pequena corrente Na$^+$ persistente
\end{itemize}

\vspace{0.3cm}

\textbf{Saída de carga (+):}
\begin{itemize}
    \item I$_{Kr}$: corrente retificadora rápida
    \item I$_{Ks}$: corrente retificadora lenta
    \item I$_{to}$: corrente transitória externa
\end{itemize}

\vspace{0.3cm}

\importante{
Equilíbrio delicado:\\
$I_{Ca,L} \approx I_{Kr} + I_{Ks}$
}

\column{0.5\textwidth}
\begin{center}
\begin{tikzpicture}[scale=0.8]
    % Voltage trace
    \draw[->] (0,0) -- (5,0) node[right, font=\tiny] {t};
    \draw[->] (0,0) -- (0,4) node[above, font=\tiny] {V$_m$};

    % AP trace
    \draw[thick, azulescuro] (0.5,0.3) -- (1,3.5) -- (1.3,2.8) -- (3,2.7) -- (4,0.3);

    % Current traces below
    \begin{scope}[yshift=-2.5cm]
        \draw[->] (0,0) -- (5,0) node[right, font=\tiny] {t};
        \draw[->] (0,-1.5) -- (0,1) node[above, font=\tiny] {I};

        % INa
        \draw[thick, verdeacido] (1,0) -- (1,-1.2) -- (1.2,-1.2) -- (1.2,0);
        \node[left, font=\tiny] at (1,-0.6) {I$_{Na}$};

        % ICaL
        \draw[thick, laranjacel] (1.2,0) -- (1.2,-0.6) -- (3.2,-0.6) -- (3.2,0);
        \node[above, font=\tiny] at (2.2,-0.4) {I$_{Ca,L}$};

        % IK
        \draw[thick, roxodna] (1.5,0) -- (1.5,0.6) -- (4.2,0.8) -- (4.2,0);
        \node[above, font=\tiny] at (3,0.7) {I$_K$};
    \end{scope}
\end{tikzpicture}
\end{center}
\end{columns>

\vspace{0.3cm}

\begin{exampleblock}{Importância Clínica}
Bloqueio de I$_{Kr}$ (drogas, mutações) $\rightarrow$ prolongamento QT $\rightarrow$ torsades de pointes
\end{exampleblock}
\end{frame}

\begin{frame}{Período Refratário Efetivo}
\definicao{Período Refratário}{
Intervalo após um potencial de ação durante o qual a célula não pode ser reexcitada (período refratário absoluto) ou requer estímulo maior (período refratário relativo).
}

\vspace{0.3cm}

\begin{columns}
\column{0.5\textwidth}
\textbf{Mecanismo:}
\begin{itemize}
    \item Inativação de canais Na$^+$
    \item Recuperação dependente de voltagem
    \item Recuperação lenta ($\sim$200 ms)
    \item Prolonga com frequência
\end{itemize}

\vspace{0.3cm}

\textbf{Significado:}
\begin{itemize}
    \item Previne reentrada
    \item Limita frequência máxima
    \item Protege contra tetania
\end{itemize}

\column{0.5\textwidth}
\begin{center}
\begin{tikzpicture}[scale=0.7]
    % AP
    \draw[->] (0,0) -- (6,0) node[right, font=\tiny] {Tempo};
    \draw[->] (0,-0.5) -- (0,3.5) node[above, font=\tiny] {V$_m$};

    % First AP
    \draw[thick, azulescuro] (0.5,0) -- (1,3) -- (1.3,2.5) -- (3,2.4) -- (3.5,0) -- (4.5,0);

    % Refractory periods
    \fill[red!20, opacity=0.5] (0.5,0) rectangle (3.5,3.5);
    \node[font=\tiny, rotate=90] at (0.3,1.5) {Absoluto};

    \fill[yellow!20, opacity=0.5] (3.5,0) rectangle (4.5,3.5);
    \node[font=\tiny, rotate=90] at (3.3,1.5) {Relativo};

    % Attempted stimuli
    \draw[->, red, thick] (2,-0.3) -- (2,0.3) node[above, font=\tiny] {\xmark};
    \draw[->, verdeacido, thick] (4,-0.3) -- (4,0.3) node[above, font=\tiny] {\checkmark};

    % Second AP (reduced)
    \draw[thick, azulescuro, dashed] (4.5,0) -- (4.8,2) -- (5,1.7) -- (6,0);
\end{tikzpicture}
\end{center>
\end{columns>

\vspace{0.3cm}

\importante{
ERP cardíaco ($\sim$200-400 ms) é muito maior que neuronal ($\sim$1-2 ms)
}
\end{frame}

\begin{frame}{Formalismo de Hodgkin-Huxley}
\begin{block}{Modelo Matemático de Canais Iônicos}
Hodgkin e Huxley (1952) descreveram canais como gates probabilísticos:

\begin{equation*}
I_{ion} = g_{ion} \cdot (V_m - E_{ion})
\end{equation*}

onde a condutância é modulada por variáveis de gating:
\begin{equation*}
g_{ion} = \bar{g}_{ion} \cdot m^p \cdot h^q
\end{equation*}

\begin{itemize}
    \item $m$: ativação (abertura)
    \item $h$: inativação (fechamento)
    \item $p, q$: expoentes de cooperatividade
\end{itemize}
\end{block}

\vspace{0.3cm}

\begin{exampleblock}{Dinâmica de Gating}
\begin{equation*}
\frac{dm}{dt} = \alpha_m(V)(1-m) - \beta_m(V)m = \frac{m_\infty(V) - m}{\tau_m(V)}
\end{equation*}
\end{exampleblock}
\end{frame>

\begin{frame}{Corrente de Sódio: I$_{Na}$}
\begin{block}{Modelo de I$_{Na}$}
\begin{equation*}
I_{Na} = g_{Na} m^3 h j (V_m - E_{Na})
\end{equation*}

\textbf{Três gates:}
\begin{itemize}
    \item $m$: ativação rápida ($\tau \sim 0.1$ ms)
    \item $h$: inativação rápida ($\tau \sim 1-5$ ms)
    \item $j$: inativação lenta ($\tau \sim 20-200$ ms)
\end{itemize}
\end{block}

\vspace{0.3cm}

\begin{columns}
\column{0.5\textwidth}
\textbf{Steady-state (exemplo):}
\begin{align*}
m_\infty &= \frac{1}{1 + e^{-(V+40)/5}} \\
h_\infty &= \frac{1}{1 + e^{(V+65)/7}}
\end{align*}

\column{0.5\textwidth}
\textbf{Time constants (exemplo):}
\begin{align*}
\tau_m &= 0.12 \text{ ms} \\
\tau_h &= \frac{1}{0.13(1+e^{-(V+10)/10})}
\end{align*}
\end{columns}

\vspace{0.3cm}

\importante{
Bloqueio de I$_{Na}$ (classe I antiarrítmicos) reduz velocidade de condução
}
\end{frame>

\begin{frame}{Correntes de Potássio e Cálcio}
\begin{columns}
\column{0.5\textwidth}
\textbf{I$_{Ca,L}$ (Corrente-L de Ca$^{2+}$):}
\begin{equation*}
I_{Ca,L} = g_{Ca} d f f_{Ca} (V - E_{Ca})
\end{equation*}
\begin{itemize}
    \item $d$: ativação
    \item $f$: inativação voltagem-dependente
    \item $f_{Ca}$: inativação Ca$^{2+}$-dependente
    \item Ativa em $\sim$-40 mV
    \item Crucial para contração
\end{itemize}

\vspace{0.3cm}

\textbf{Bloqueadores:}
\begin{itemize}
    \item Dihidropiridinas (nifedipina)
    \item Verapamil, diltiazem
\end{itemize}

\column{0.5\textwidth}
\textbf{I$_{Kr}$ (Rapid delayed rectifier):}
\begin{equation*}
I_{Kr} = g_{Kr} x_{Kr1} x_{Kr2} (V - E_K)
\end{equation*}
\begin{itemize}
    \item hERG channel (Kv11.1)
    \item Ativação lenta
    \item Retificação interna
    \item Repolarização Fase 3
\end{itemize}

\vspace{0.3cm}

\textbf{I$_{Ks}$ (Slow delayed rectifier):}
\begin{equation*}
I_{Ks} = g_{Ks} x_s^2 (V - E_K)
\end{equation*}
\begin{itemize}
    \item KCNQ1/KCNE1
    \item Ativação muito lenta
    \item Importante em frequências altas
\end{itemize}
\end{columns}
\end{frame>

\begin{frame}[fragile]{Simulação: Potencial de Ação Simples}
\begin{lstlisting}[language=Python, caption=Modelo simplificado de PA cardíaco, basicstyle=\ttfamily\footnotesize]
import numpy as np
import matplotlib.pyplot as plt

def simple_cardiac_ap(t, V, gates):
    """Modelo simplificado de PA cardiaco"""
    m, h, d, f, x = gates

    # Parametros
    gNa, gCa, gK = 23, 0.09, 0.282
    ENa, ECa, EK = 50, 60, -85

    # Correntes ionicas
    INa = gNa * m**3 * h * (V - ENa)
    ICa = gCa * d * f * (V - ECa)
    IK = gK * x * (V - EK)

    # Capacitancia
    Cm = 1.0  # uF/cm^2

    # dV/dt
    dVdt = -(INa + ICa + IK) / Cm

    # Gating variables (simplified)
    m_inf = 1/(1 + np.exp(-(V+40)/5))
    h_inf = 1/(1 + np.exp((V+65)/7))
    d_inf = 1/(1 + np.exp(-(V+10)/6.24))
    f_inf = 1/(1 + np.exp((V+32)/8))
    x_inf = 1/(1 + np.exp(-(V+20)/15))

    tau_m, tau_h = 0.1, 2.0
    tau_d, tau_f, tau_x = 3.0, 20.0, 100.0

    dmdt = (m_inf - m) / tau_m
    dhdt = (h_inf - h) / tau_h
    dddt = (d_inf - d) / tau_d
    dfdt = (f_inf - f) / tau_f
    dxdt = (x_inf - x) / tau_x

    return dVdt, [dmdt, dhdt, dddt, dfdt, dxdt]
\end{lstlisting}
\end{frame>

% ============================================================================
% SEÇÃO 4: MODELOS DE CARDIOMIÓCITOS
% ============================================================================
\section{Modelos de Cardiomiócitos}

\begin{frame}{História dos Modelos Cardíacos}
\begin{block}{Evolução Temporal}
\begin{center}
\begin{tabular}{lll}
\toprule
\textbf{Ano} & \textbf{Modelo} & \textbf{Características} \\
\midrule
1962 & Noble & Primeira célula cardíaca (Purkinje) \\
     &       & 4 correntes, 6 variáveis \\
\midrule
1977 & Beeler-Reuter & Ventricular mamífero \\
     &               & 4 correntes, 8 variáveis \\
\midrule
1991 & Luo-Rudy I (LR1) & Dinâmica de Ca$^{2+}$ \\
     &                   & 6 correntes \\
\midrule
1994 & Luo-Rudy II (LR2) & Ca$^{2+}$ intracelular detalhado \\
\midrule
2004 & ten Tusscher & Ventricular humano \\
     &              & Epi, M, Endo cells \\
\midrule
2011 & O'Hara-Rudy & Ventricular humano \\
     &             & Baseado em dados experimentais \\
\bottomrule
\end{tabular}
\end{center}
\end{block}

\importante{
Modelos mais recentes: $>$15 correntes, $>$40 variáveis de estado
}
\end{frame>

\begin{frame}{Modelo de Noble (1962)}
\exemplo{Primeiro Modelo Cardíaco}{
Denis Noble adaptou o modelo de Hodgkin-Huxley para fibras de Purkinje cardíacas.
}

\vspace{0.3cm}

\begin{columns}
\column{0.5\textwidth}
\textbf{Correntes:}
\begin{itemize}
    \item I$_{Na}$: corrente de sódio
    \item I$_{K}$: corrente de potássio
    \item I$_{leak}$: corrente de vazamento
    \item I$_{an}$: corrente aniônica
\end{itemize}

\vspace{0.3cm}

\textbf{Equação de Membrana:}
\begin{equation*}
C_m \frac{dV}{dt} = -I_{Na} - I_K - I_{leak} - I_{an}
\end{equation*}

\column{0.5\textwidth}
\textbf{Realizações:}
\begin{itemize}
    \item Reproduziu PA de Purkinje
    \item Automaticidade espontânea
    \item Resposta a estímulos
    \item Base para modelos futuros
\end{itemize}

\vspace{0.3cm}

\begin{alertblock}{Limitações}
\begin{itemize}
    \item Sem Ca$^{2+}$ intracelular
    \item Sem compartimentos
    \item Cinética simplificada
\end{itemize}
\end{alertblock}
\end{columns}
\end{frame>

\begin{frame}{Modelo de Luo-Rudy (LR1, 1991)}
\begin{block}{Avanços do LR1}
\begin{itemize}
    \item Primeiro modelo ventricular com dinâmica de Ca$^{2+}$ detalhada
    \item Retículo sarcoplasmático (SR)
    \item Bomba Na$^+$/K$^+$-ATPase
    \item Trocador Na$^+$/Ca$^{2+}$ (NCX)
    \item Acoplamento excitação-contração
\end{itemize}
\end{block}

\vspace{0.3cm}

\begin{columns}
\column{0.5\textwidth}
\textbf{Correntes (6):}
\begin{itemize}
    \item I$_{Na}$, I$_{Ca,L}$
    \item I$_{K}$, I$_{K1}$
    \item I$_{Kp}$
    \item I$_{b}$ (background)
\end{itemize>

\textbf{Bombas e Trocadores:}
\begin{itemize}
    \item Na$^+$/K$^+$ pump
    \item Na$^+$/Ca$^{2+}$ exchanger
    \item Ca$^{2+}$ pump (sarcolemma)
\end{itemize}

\column{0.5\textwidth}
\textbf{Compartimentos:}
\begin{itemize}
    \item Citosol
    \item SR (uptake/release)
\end{itemize}

\vspace{0.3cm}

\begin{center}
\begin{tikzpicture}[scale=0.6]
    % Cell membrane
    \draw[thick] (0,0) rectangle (4,3);
    \node[above] at (2,3) {\tiny Sarcolema};

    % SR
    \draw[thick, roxodna] (0.5,0.5) rectangle (3.5,2.5);
    \node at (2,1.5) {\tiny SR};

    % Ca fluxes
    \draw[->, thick, red] (0.2,2) -- node[left, font=\tiny] {I$_{Ca,L}$} (0.2,1);
    \draw[->, thick, verdeacido] (2,0.5) -- node[below, font=\tiny] {Uptake} (2,0.2);
    \draw[->, thick, azulclaro] (2,2.5) -- node[above, font=\tiny] {Release} (2,2.8);
\end{tikzpicture}
\end{center>
\end{columns>
\end{frame>

\begin{frame}{Modelo ten Tusscher (2004)}
\begin{block}{Modelo Ventricular Humano Detalhado}
Desenvolvido a partir de dados experimentais humanos
\end{block}

\vspace{0.3cm}

\begin{columns}
\column{0.5\textwidth}
\textbf{Três Tipos Celulares:}
\begin{itemize}
    \item \textbf{Endocárdio:} PA mais curto
    \item \textbf{M cells:} PA mais longo (mid-myocardium)
    \item \textbf{Epicárdio:} I$_{to}$ proeminente
\end{itemize}

\vspace{0.3cm}

\textbf{Correntes (12):}
\begin{itemize}
    \item I$_{Na}$, I$_{to}$, I$_{Kr}$, I$_{Ks}$
    \item I$_{K1}$, I$_{Ca,L}$, I$_{NaCa}$
    \item I$_{NaK}$, I$_{pCa}$, I$_{pK}$
    \item I$_{bNa}$, I$_{bCa}$
\end{itemize}

\column{0.5\textwidth}
\textbf{Aplicações:}
\begin{itemize}
    \item Dispersão transmural da repolarização
    \item Efeitos pró-arrítmicos de drogas
    \item Síndrome do QT longo
    \item Fibrilação ventricular
\end{itemize}

\vspace{0.3cm}

\begin{exampleblock}{Heterogeneidade}
A diferença entre tipos celulares cria gradiente de repolarização que gera a onda T do ECG
\end{exampleblock>
\end{columns>

\vspace{0.3cm}

\importante{
Modelo amplamente usado: $>$500 citações, base para estudos de arritmias
}
\end{frame>

\begin{frame}{Comparação: Modelos Atriais vs Ventriculares}
\begin{center}
\begin{tabular}{lll}
\toprule
\textbf{Característica} & \textbf{Atrial} & \textbf{Ventricular} \\
\midrule
Duração PA & 150-300 ms & 200-400 ms \\
Platô & Menos pronunciado & Pronunciado \\
I$_{to}$ & Grande & Variável \\
I$_{K,ACh}$ & Presente & Ausente \\
I$_{Kur}$ & Presente (atrial) & Ausente \\
APD adaptação & Rápida & Lenta \\
Velocidade condução & 0.5-1 m/s & 0.3-0.5 m/s \\
\bottomrule
\end{tabular}
\end{center>

\vspace{0.3cm}

\begin{columns}
\column{0.5\textwidth}
\textbf{Modelos Atriais:}
\begin{itemize}
    \item Courtemanche et al. (1998)
    \item Nygren et al. (1998)
    \item Grandi et al. (2011)
\end{itemize}

\column{0.5\textwidth}
\begin{alertblock}{Diferenças Importantes}
Modelos atriais essenciais para estudar fibrilação atrial, a arritmia cardíaca mais comum
\end{alertblock}
\end{columns>
\end{frame>

\begin{frame}{Complexidade Computacional}
\begin{block}{Trade-off: Realismo vs Custo}
\begin{center}
\begin{tikzpicture}[scale=0.8]
    % Axes
    \draw[->] (0,0) -- (8,0) node[right] {Complexidade};
    \draw[->] (0,0) -- (0,5) node[above] {Realismo/Custo};

    % Curve for realism
    \draw[thick, verdeacido, domain=0:7, samples=100]
        plot (\x, {3.5*(1-exp(-\x/2))}) node[right] {\tiny Realismo};

    % Curve for computational cost
    \draw[thick, red, domain=0:7, samples=100]
        plot (\x, {0.3*\x}) node[right] {\tiny Custo};

    % Model points
    \node[circle, fill=azulclaro, minimum size=0.3cm] at (1.5,1.5) {};
    \node[below, font=\tiny] at (1.5,1) {Noble};

    \node[circle, fill=azulclaro, minimum size=0.3cm] at (3,2.5) {};
    \node[below, font=\tiny] at (3,2) {LR1};

    \node[circle, fill=azulclaro, minimum size=0.3cm] at (5,3.2) {};
    \node[below, font=\tiny] at (5,2.7) {ten Tusscher};

    \node[circle, fill=azulclaro, minimum size=0.3cm] at (6.5,3.4) {};
    \node[below, font=\tiny] at (6.5,2.9) {O'Hara-Rudy};
\end{tikzpicture}
\end{center>
\end{block>

\importante{
Para simulações de tecido 3D: modelos simplificados (FitzHugh-Nagumo, Aliev-Panfilov)
}
\end{frame>

\begin{frame}[fragile]{Implementação: Modelo Simplificado}
\begin{lstlisting}[language=Python, caption=Modelo cardíaco simplificado, basicstyle=\ttfamily\footnotesize]
import numpy as np
from scipy.integrate import odeint

def cardiac_model(y, t, I_stim):
    """Modelo simplificado baseado em LR1"""
    V, m, h, j, d, f, X = y

    # Parametros
    gNa, gCa, gK, gK1 = 23, 0.09, 0.282, 0.6047
    ENa, ECa, EK = 54.4, 65, -77

    # Gating steady-states e time constants
    m_inf = 1/(1 + np.exp(-(V+47.13)/10.3))
    h_inf = 1/(1 + np.exp((V+78.5)/6.22))
    # ... (outros gates)

    # Correntes
    INa = gNa * m**3 * h * j * (V - ENa)
    ICa = gCa * d * f * (V - ECa)
    IK = gK * X * (V - EK)
    IK1 = gK1 * (V - EK) / (1 + np.exp(0.07*(V+80)))

    # dV/dt
    Cm = 1.0
    dVdt = -(INa + ICa + IK + IK1 - I_stim) / Cm

    # Gating kinetics
    dmdt = (m_inf - m) / tau_m
    dhdt = (h_inf - h) / tau_h
    # ... (outros gates)

    return [dVdt, dmdt, dhdt, djdt, dddt, dfdt, dXdt]

# Simulacao
t = np.linspace(0, 500, 5000)
y0 = [-85, 0.01, 0.99, 0.99, 0.001, 0.99, 0.01]
I_stim = lambda t: 50 if 5 < t < 6 else 0

sol = odeint(cardiac_model, y0, t, args=(I_stim,))
\end{lstlisting}
\end{frame>

% ============================================================================
% SEÇÃO 5: ACOPLAMENTO EXCITAÇÃO-CONTRAÇÃO
% ============================================================================
\section{Acoplamento Excitação-Contração}

\begin{frame}{Acoplamento Excitação-Contração (EC)}
\definicao{Acoplamento EC}{
Processo pelo qual o potencial de ação elétrico desencadeia a contração mecânica do cardiomiócito através da regulação do Ca$^{2+}$ intracelular.
}

\vspace{0.3cm}

\begin{block}{Sequência de Eventos}
\begin{enumerate}
    \item PA despolariza membrana
    \item Abertura de canais Ca$^{2+}$ tipo-L (DHPR)
    \item Entrada de Ca$^{2+}$ no citosol (I$_{Ca,L}$)
    \item \textbf{Ca$^{2+}$-induced Ca$^{2+}$ release (CICR)}
    \item Receptores de rianodina (RyR) no SR liberam Ca$^{2+}$
    \item $[Ca^{2+}]_i$ aumenta de $\sim$0.1 $\mu$M para $\sim$1-10 $\mu$M
    \item Ca$^{2+}$ liga-se à troponina C
    \item Contração dos filamentos de actina-miosina
    \item Recaptação de Ca$^{2+}$ pelo SR (SERCA)
    \item Relaxamento
\end{enumerate}
\end{block>
\end{frame>

\begin{frame}{CICR: Liberação de Ca$^{2+}$ Induzida por Ca$^{2+}$}
\begin{columns}
\column{0.5\textwidth}
\textbf{Mecanismo:}
\begin{itemize}
    \item Pequena entrada Ca$^{2+}$ via DHPR
    \item Ativa receptores RyR2 no SR
    \item Liberação massiva de Ca$^{2+}$
    \item Amplificação: 1:10 ou maior
    \item "Calcium sparks"
\end{itemize}

\vspace{0.3cm}

\textbf{Regulação:}
\begin{itemize}
    \item [Ca$^{2+}$]$_{SR}$: carga do SR
    \item Fosforilação de RyR
    \item Calmodulina
    \item ATP, Mg$^{2+}$
\end{itemize}

\column{0.5\textwidth}
\begin{center}
\begin{tikzpicture}[scale=0.7]
    % T-tubule and SR
    \draw[thick] (0,0) -- (4,0) node[right, font=\tiny] {T-tubule};
    \draw[thick, roxodna] (0,-2) -- (4,-2) node[right, font=\tiny] {SR};

    % DHPR (L-type Ca channel)
    \draw[thick, fill=laranjacel!50] (1.5,-0.3) rectangle (1.8,0.3);
    \node[above, font=\tiny] at (1.65,0.5) {DHPR};

    % RyR
    \draw[thick, fill=verdeacido!50] (1.5,-2.3) rectangle (1.8,-1.7);
    \node[below, font=\tiny] at (1.65,-2.5) {RyR};

    % Ca influx
    \draw[->, ultra thick, red] (1.65,0.3) -- (1.65,-0.8) node[midway, right, font=\tiny] {Ca$^{2+}$};

    % CICR
    \draw[->, ultra thick, red] (1.65,-1) to[out=-90, in=90] (1.65,-1.7);
    \node[right, font=\tiny] at (1.9,-1.35) {Trigger};

    % Ca release from SR
    \draw[->, ultra thick, red] (1.65,-2.3) -- (1.65,-3) node[below, font=\tiny] {Ca$^{2+}$ release};

    % Amplification
    \foreach \x in {0.8, 1.1, 1.4, 2, 2.3, 2.6} {
        \draw[->, red, thick] (\x,-2.3) -- (\x,-2.8);
    }
\end{tikzpicture}
\end{center>
\end{columns>

\vspace{0.3cm}

\importante{
Falha do CICR: causa principal de insuficiência cardíaca sistólica
}
\end{frame>

\begin{frame}{Dinâmica de Ca$^{2+}$ no Retículo Sarcoplasmático}
\begin{block}{Compartimentos de Ca$^{2+}$}
\begin{itemize}
    \item \textbf{Citosol:} [Ca$^{2+}$]$_i$ $\sim$0.1 $\mu$M (diástole) $\rightarrow$ 1-10 $\mu$M (sístole)
    \item \textbf{SR:} [Ca$^{2+}$]$_{SR}$ $\sim$1 mM (reservatório)
    \item \textbf{Subsarcolemmal:} [Ca$^{2+}$]$_{ss}$ (microdomínio)
\end{itemize}
\end{block>

\vspace{0.3cm}

\begin{columns}
\column{0.5\textwidth}
\textbf{Fluxos:}
\begin{itemize}
    \item \textbf{Uptake (SERCA):}\\
    $J_{up} = V_{max,up} \frac{[Ca^{2+}]_i^2}{K_m^2 + [Ca^{2+}]_i^2}$

    \item \textbf{Release (RyR):}\\
    $J_{rel} = k_{rel} [Ca^{2+}]_{SR} \cdot O_{RyR}$

    \item \textbf{Leak:}\\
    $J_{leak} = k_{leak}([Ca^{2+}]_{SR} - [Ca^{2+}]_i)$
\end{itemize}

\column{0.5\textwidth}
\textbf{Equação de Balanço:}
\begin{equation*}
\frac{d[Ca^{2+}]_i}{dt} = J_{rel} - J_{up} + J_{leak} + \frac{I_{Ca,L}}{2FV_{cell}}
\end{equation*}

\vspace{0.3cm}

\textbf{Para o SR:}
\begin{equation*}
\frac{d[Ca^{2+}]_{SR}}{dt} = J_{up} - J_{rel} - J_{leak}
\end{equation*}
\end{columns>

\vspace{0.3cm}

\begin{exampleblock}{Buffers}
Ca$^{2+}$ é tamponado por troponina C, calmodulina, calsequestrin (SR)
\end{exampleblock>
\end{frame>

\begin{frame}{Geração de Força: Cross-Bridge Cycling}
\begin{columns}
\column{0.5\textwidth}
\textbf{Ciclo de Pontes Cruzadas:}
\begin{enumerate}
    \item Ca$^{2+}$ liga troponina C
    \item Conformação da tropomiosina
    \item Exposição de sítios de ligação
    \item Miosina liga actina
    \item Power stroke (ADP liberado)
    \item ATP liga miosina
    \item Desligamento
    \item Hidrólise ATP (recarga)
\end{enumerate}

\vspace{0.3cm}

\textbf{Força:}
\begin{equation*}
F = F_{max} \cdot f([Ca^{2+}]_i) \cdot g(L)
\end{equation*}

\column{0.5\textwidth}
\begin{center}
\begin{tikzpicture}[scale=0.6]
    % Thin filament (actin)
    \draw[thick, verdeacido] (0,2) -- (5,2);
    \node[above, font=\tiny] at (2.5,2) {Actina};

    % Thick filament (myosin)
    \draw[thick, roxodna] (0,0) -- (5,0);
    \node[below, font=\tiny] at (2.5,0) {Miosina};

    % Myosin heads
    \foreach \x in {1, 2, 3, 4} {
        \draw[thick] (\x,0) -- (\x,0.8);
        \draw[thick, fill=laranjacel!50] (\x,0.8) circle (0.15);
    }

    % Ca2+ and Troponin
    \draw[fill=red, circle] (2.5,2.5) circle (0.1) node[above, font=\tiny] {Ca$^{2+}$};
    \draw[thick, azulclaro] (2.3,2) rectangle (2.7,2.2);
    \node[right, font=\tiny] at (2.8,2.1) {TnC};

    % Power stroke
    \draw[->, ultra thick, red] (2,0.8) -- (2.3,1.7);
    \node[right, font=\tiny] at (2.3,1.2) {Power stroke};

    % Sarcomere shortening
    \draw[<->, thick] (0,-0.5) -- (5,-0.5);
    \node[below, font=\tiny] at (2.5,-0.5) {Encurtamento};
\end{tikzpicture}
\end{center>

\vspace{0.3cm}

\textbf{Dependência de Ca$^{2+}$:}
\begin{equation*}
f([Ca^{2+}]) = \frac{[Ca^{2+}]_i^n}{K_d^n + [Ca^{2+}]_i^n}
\end{equation*}
Hill coefficient $n \approx 4-6$
\end{columns>
\end{frame>

\begin{frame}{Relação Comprimento-Tensão (Frank-Starling)}
\definicao{Lei de Frank-Starling}{
A força de contração do coração é proporcional ao comprimento inicial das fibras musculares (pré-carga), dentro de limites fisiológicos.
}

\vspace{0.3cm}

\begin{columns}
\column{0.5\textwidth}
\textbf{Mecanismo Molecular:}
\begin{itemize}
    \item Sobreposição actina-miosina ótima
    \item Sensibilidade ao Ca$^{2+}$ aumenta com comprimento
    \item Alinhamento de filamentos
    \item Recrutamento de pontes cruzadas
\end{itemize}

\vspace{0.3cm}

\textbf{Significado Fisiológico:}
\begin{itemize}
    \item Ajuste do débito cardíaco
    \item Compensação de alterações de pré-carga
    \item Equilíbrio entre ventrículos
\end{itemize>

\column{0.5\textwidth}
\begin{center}
\begin{tikzpicture}[scale=0.8]
    % Force-length curve
    \draw[->] (0,0) -- (5,0) node[right, font=\tiny] {Comprimento (L/L$_0$)};
    \draw[->] (0,0) -- (0,4) node[above, font=\tiny] {Força};

    % Ascending limb
    \draw[thick, azulescuro, domain=0.8:2, samples=50]
        plot (\x, {2.5*(\x-0.8)});

    % Plateau
    \draw[thick, azulescuro] (2,3) -- (2.8,3);

    % Descending limb
    \draw[thick, azulescuro, domain=2.8:4.5, samples=50]
        plot (\x, {3 - 0.8*(\x-2.8)});

    % Optimal length
    \draw[dashed, red] (2.4,0) -- (2.4,3);
    \node[below, font=\tiny] at (2.4,0) {L$_0$};

    % Regions
    \node[font=\tiny] at (1.5,1) {Subótimo};
    \node[font=\tiny] at (2.4,3.5) {Ótimo};
    \node[font=\tiny] at (3.7,1) {Superestirado};
\end{tikzpicture}
\end{center>
\end{columns>

\vspace{0.3cm}

\importante{
Comprimento ótimo do sarcômero: 2.0-2.2 $\mu$m (máxima sobreposição)
}
\end{frame>

\begin{frame}[fragile]{Modelo de Contração}
\begin{lstlisting}[language=Python, caption=Modelo simplificado de contração, basicstyle=\ttfamily\footnotesize]
def contraction_model(Ca_i, SL):
    """
    Modelo de geracao de forca
    Ca_i: concentracao de calcio intracelular (uM)
    SL: sarcomere length (um)
    """
    # Parametros
    Ca_50 = 1.0  # uM, [Ca] para 50% ativacao
    n_Hill = 4   # Coeficiente de Hill
    F_max = 100  # kPa, forca maxima
    L_0 = 2.0    # um, comprimento otimo

    # Dependencia de Ca2+
    Ca_activation = Ca_i**n_Hill / (Ca_50**n_Hill + Ca_i**n_Hill)

    # Dependencia de comprimento (Frank-Starling)
    if SL < 1.6:
        length_dep = 0
    elif SL < L_0:
        length_dep = (SL - 1.6) / (L_0 - 1.6)
    elif SL < 2.4:
        length_dep = 1.0
    else:
        length_dep = 1.0 - 0.5*(SL - 2.4)

    # Forca total
    Force = F_max * Ca_activation * length_dep

    return Force

# Exemplo
import numpy as np
import matplotlib.pyplot as plt

Ca_range = np.linspace(0.1, 10, 100)
SL_values = [1.8, 2.0, 2.2]

plt.figure(figsize=(10, 6))
for SL in SL_values:
    F = [contraction_model(ca, SL) for ca in Ca_range]
    plt.plot(Ca_range, F, linewidth=2, label=f'SL = {SL} um')

plt.xlabel('[Ca2+]i (uM)')
plt.ylabel('Forca (kPa)')
plt.legend()
plt.grid(True)
\end{lstlisting>
\end{frame>

% ============================================================================
% SEÇÃO 6: PROPAGAÇÃO E ARRITMIAS
% ============================================================================
\section{Propagação e Arritmias}

\begin{frame}{Propagação da Onda Elétrica}
\begin{block}{Teoria do Cabo (Cable Theory)}
Propagação do potencial em tecido cardíaco:
\begin{equation*}
\lambda^2 \frac{\partial^2 V}{\partial x^2} = V - V_{rest} + \tau \frac{\partial V}{\partial t}
\end{equation*}

onde:
\begin{itemize}
    \item $\lambda$: constante de comprimento espacial
    \item $\tau$: constante de tempo
    \item $V$: potencial transmembrana
\end{itemize}
\end{block>

\vspace{0.3cm}

\begin{columns}
\column{0.5\textwidth}
\textbf{Equação de Reação-Difusão:}
\begin{equation*}
\frac{\partial V}{\partial t} = D\nabla^2 V - \frac{I_{ion}}{C_m}
\end{equation*}

\begin{itemize}
    \item $D$: coeficiente de difusão
    \item $\nabla^2$: operador laplaciano
    \item $I_{ion}$: correntes iônicas
\end{itemize>

\column{0.5\textwidth}
\textbf{Velocidade de Condução:}
\begin{equation*}
\theta = \lambda \sqrt{\frac{2}{\tau}}
\end{equation*}

\vspace{0.3cm}

\textbf{Fatores que Afetam $\theta$:}
\begin{itemize}
    \item Densidade de Na$^+$ channels
    \item Resistência gap junctions
    \item Geometria do tecido
    \item Refratariedade
    \item Isquemia, fibrose
\end{itemize>
\end{columns>
\end{frame>

\begin{frame}{Mecanismos de Arritmias}
\begin{block}{Classificação}
\begin{enumerate}
    \item \textbf{Anormalidades de Formação do Impulso:}
    \begin{itemize}
        \item Automaticidade anormal
        \item Atividade deflagrada (DADs, EADs)
    \end{itemize}

    \item \textbf{Anormalidades de Condução:}
    \begin{itemize}
        \item Bloqueios (AV, ramo)
        \item Condução lenta
    \end{itemize}

    \item \textbf{Reentrada:}
    \begin{itemize}
        \item Circuito anatômico
        \item Circuito funcional (spiral waves)
    \end{itemize>
\end{enumerate}
\end{block>

\vspace{0.3cm}

\importante{
Reentrada é o mecanismo mais comum de taquiarritmias sustentadas (FA, FV, TV)
}
\end{frame>

\begin{frame}{Reentrada e Ondas Espirais}
\begin{columns}
\column{0.5\textwidth}
\textbf{Critérios para Reentrada:}
\begin{enumerate}
    \item Circuito de condução
    \item Bloqueio unidirecional
    \item Condução lenta
    \item ERP $<$ comprimento do circuito
\end{enumerate}

\vspace{0.3cm}

\textbf{Comprimento de Onda ($\lambda$):}
\begin{equation*}
\lambda = \theta \times ERP
\end{equation*}

\vspace{0.3cm}

\begin{alertblock}{Condição de Reentrada}
\begin{equation*}
\text{Perímetro do circuito} > \lambda
\end{equation*}
\end{alertblock>

\column{0.5\textwidth}
\begin{center}
\begin{tikzpicture}[scale=0.8]
    % Reentry circuit
    \draw[ultra thick, azulescuro] (0,0) circle (1.5);

    % Obstacle
    \draw[fill=gray!50] (0,0) circle (0.5);
    \node at (0,0) {\tiny Obstáculo};

    % Wave propagation (arrows)
    \foreach \angle in {0, 45, 90, 135, 180, 225, 270, 315} {
        \draw[->, very thick, red]
            ({1.5*cos(\angle)}, {1.5*sin(\angle)})
            arc (\angle:\angle+40:1.5);
    }

    % Refractory region
    \draw[thick, fill=yellow!30, opacity=0.5]
        (1.5,0) arc (0:90:1.5) -- (0.5,1.5) arc (90:0:0.5) -- cycle;
    \node[font=\tiny] at (1.2,0.8) {Refratário};

    % Excitable region
    \draw[thick, fill=verdeacido!30, opacity=0.5]
        (0,1.5) arc (90:180:1.5) -- (-0.5,0) arc (180:90:0.5) -- cycle;
    \node[font=\tiny] at (-1.2,0.8) {Excitável};
\end{tikzpicture}
\end{center>

\textbf{Ondas Espirais:}
\begin{itemize}
    \item Rotor: centro de rotação
    \item Frequência alta
    \item Sustentadas
    \item Fibrilação
\end{itemize>
\end{columns>
\end{frame>

\begin{frame}{Fibrilação Atrial (FA)}
\begin{block}{Epidemiologia e Impacto}
\begin{itemize}
    \item Arritmia sustentada mais comum
    \item Prevalência: 1-2\% da população geral
    \item Aumenta com idade ($>$10\% em $>$80 anos)
    \item Risco de AVC aumentado 5x
    \item Insuficiência cardíaca
\end{itemize}
\end{block>

\vspace{0.3cm}

\begin{columns}
\column{0.5\textwidth}
\textbf{Mecanismos:}
\begin{itemize}
    \item \textbf{Trigger:} focos ectópicos (veias pulmonares)
    \item \textbf{Substrato:} remodelamento atrial
    \item Multiple wavelets
    \item Rotores
    \item Fibrose
\end{itemize}

\vspace{0.3cm}

\textbf{Remodelamento Elétrico:}
\begin{itemize}
    \item $\downarrow$ I$_{Ca,L}$
    \item $\uparrow$ I$_{K1}$
    \item APD encurtado
    \item ERP reduzido
    \item "AF begets AF"
\end{itemize>

\column{0.5\textwidth}
\textbf{Tratamento:}
\begin{itemize}
    \item \textbf{Controle de frequência:}\\
    Beta-bloqueadores, diltiazem

    \item \textbf{Controle de ritmo:}\\
    Antiarrítmicos (classe I, III)\\
    Cardioversão

    \item \textbf{Ablação por cateter:}\\
    Isolamento de veias pulmonares

    \item \textbf{Anticoagulação:}\\
    Prevenção de AVC
\end{itemize>

\vspace{0.3cm}

\begin{exampleblock}{Modelagem}
Modelos computacionais de FA ajudam a:
\begin{itemize}
    \item Identificar alvos de ablação
    \item Prever resposta a drogas
    \item Personalizar tratamento
\end{itemize>
\end{exampleblock>
\end{columns>
\end{frame>

\begin{frame}{Fibrilação Ventricular (FV)}
\begin{block}{Emergência Médica}
\begin{itemize}
    \item Causa mais comum de morte súbita cardíaca
    \item Ritmo caótico, sem débito cardíaco efetivo
    \item Fatal em minutos sem tratamento
    \item Desfibrilação: único tratamento efetivo
\end{itemize}
\end{block>

\vspace{0.3cm}

\begin{columns}
\column{0.5\textwidth}
\textbf{Mecanismos:}
\begin{itemize}
    \item Múltiplas ondas espirais
    \item Break-up de spiral waves
    \item Dispersão de refratariedade
    \item Isquemia aguda
    \item Cicatriz pós-infarto
\end{itemize>

\vspace{0.3cm}

\textbf{Fatores Precipitantes:}
\begin{itemize}
    \item Infarto agudo do miocárdio
    \item Cardiomiopatia
    \item Canalopatias (Brugada, QTL)
    \item Drogas pró-arrítmicas
    \item Distúrbios eletrolíticos
\end{itemize>

\column{0.5\textwidth}
\begin{center}
\begin{tikzpicture}[scale=0.7]
    % Chaotic activation
    \draw[thick] (0,0) rectangle (4,4);

    % Multiple spirals
    \foreach \x/\y in {1/1, 2.5/3, 3.5/1.5} {
        \draw[ultra thick, red, ->] (\x,\y)
            .. controls ++(0.3,0.3) and ++(-0.3,0.3) .. ++( 0,0.5)
            .. controls ++(0.3,0) and ++(0.3,0.3) .. ++(0.5,0)
            .. controls ++(0,-0.3) and ++(0.3,-0.3) .. ++(0,-0.5);
    }

    \node[below, font=\tiny] at (2,0) {FV: Múltiplas ondas};

    % ECG trace
    \begin{scope}[yshift=-2cm, xshift=0.5cm]
        \draw[->] (0,0) -- (3.5,0) node[right, font=\tiny] {t};
        \draw[->] (0,-0.5) -- (0,0.5);

        % Chaotic ECG
        \draw[thick, azulescuro, domain=0:3, samples=100, smooth]
            plot (\x, {0.2*sin(40*\x r) + 0.1*rand});

        \node[below, font=\tiny] at (1.75,-0.7) {ECG em FV};
    \end{scope>
\end{tikzpicture}
\end{center>

\vspace{0.3cm}

\textbf{Desfibrilação:}
\begin{itemize}
    \item Choque elétrico 200-360 J
    \item Despolarização global
    \item Reset do sistema
    \item Permite retomada do ritmo sinusal
\end{itemize>
\end{columns>
\end{frame>

\begin{frame}{Síndrome do QT Longo (LQTS)}
\begin{block}{Definição}
Prolongamento do intervalo QT no ECG ($>$440 ms em homens, $>$460 ms em mulheres), associado a risco aumentado de torsades de pointes e morte súbita.
\end{block>

\vspace{0.3cm}

\begin{columns}
\column{0.5\textwidth}
\textbf{Tipos Genéticos:}
\begin{center}
\begin{tabular}{lll}
\toprule
\textbf{Tipo} & \textbf{Gene} & \textbf{Proteína} \\
\midrule
LQT1 & KCNQ1 & I$_{Ks}$ $\downarrow$ \\
LQT2 & KCNH2 & I$_{Kr}$ $\downarrow$ \\
LQT3 & SCN5A & I$_{Na}$ $\uparrow$ \\
\bottomrule
\end{tabular}
\end{center>

\vspace{0.3cm}

\textbf{LQT Adquirido:}
\begin{itemize}
    \item Drogas (antiarrítmicos, antibióticos)
    \item Distúrbios eletrolíticos (hipocalemia)
    \item Bradicardia
\end{itemize>

\column{0.5\textwidth}
\textbf{Mecanismo:}
\begin{itemize}
    \item APD prolongado
    \item Dispersão de repolarização
    \item Early afterdepolarizations (EADs)
    \item Torsades de pointes
\end{itemize>

\vspace{0.3cm}

\begin{center}
\begin{tikzpicture}[scale=0.6]
    % Normal AP
    \draw[thick, verdeacido] (0,0) -- (0.5,0) -- (0.7,2.5) -- (1,2.3) -- (2.5,2.2) -- (3,0) -- (4,0);
    \node[left, font=\tiny] at (0,1.2) {Normal};

    % LQTS AP with EAD
    \begin{scope}[yshift=-1.5cm]
        \draw[thick, red] (0,0) -- (0.5,0) -- (0.7,2.5) -- (1,2.3) -- (3,2.2);
        % EAD
        \draw[thick, red] (2,2.2) -- (2.2,2.7) -- (2.4,2.2);
        \draw[thick, red] (3,2.2) -- (4.5,0) -- (5.5,0);

        \node[left, font=\tiny] at (0,1.2) {LQTS};
        \draw[<->, azulclaro] (0.7,-0.5) -- (4.5,-0.5);
        \node[below, font=\tiny] at (2.6,-0.5) {QT prolongado};
    \end{scope>
\end{tikzpicture}
\end{center>

\textbf{Tratamento:}
\begin{itemize}
    \item Beta-bloqueadores
    \item Evitar drogas que prolongam QT
    \item CDI (cardiodesfibrilador implantável)
\end{itemize>
\end{columns>
\end{frame>

\begin{frame}{Efeitos de Drogas nos Canais Iônicos}
\begin{block}{Classes de Antiarrítmicos (Classificação de Vaughan Williams)}
\begin{center}
\begin{tabular}{llll}
\toprule
\textbf{Classe} & \textbf{Alvo} & \textbf{Efeito} & \textbf{Exemplos} \\
\midrule
I & I$_{Na}$ & $\downarrow$ condução & Quinidina, lidocaína \\
II & $\beta$-receptores & $\downarrow$ automaticidade & Propranolol \\
III & I$_K$ & $\uparrow$ APD, ERP & Amiodarona, sotalol \\
IV & I$_{Ca,L}$ & $\downarrow$ nó AV & Verapamil, diltiazem \\
\bottomrule
\end{tabular}
\end{center>
\end{block>

\vspace{0.3cm}

\begin{columns}
\column{0.5\textwidth}
\begin{exampleblock}{Efeitos Desejados}
\begin{itemize}
    \item Terminação de arritmias
    \item Prevenção de recorrência
    \item Controle de frequência
    \item Melhora sintomática
\end{itemize}
\end{exampleblock>

\column{0.5\textwidth}
\begin{alertblock}{Efeitos Adversos}
\begin{itemize}
    \item Pró-arritmia (paradoxal)
    \item Bradicardia
    \item Bloqueio AV
    \item Prolongamento QT
    \item Insuficiência cardíaca
\end{itemize}
\end{alertblock>
\end{columns>

\vspace{0.3cm}

\importante{
Modelos computacionais são usados para prever efeitos cardiotóxicos de novas drogas (teste in silico)
}
\end{frame>

% ============================================================================
% SEÇÃO 7: EXERCÍCIOS
% ============================================================================
\section{Exercícios}

\begin{frame}{Exercícios - Parte 1}
\begin{block}{Questão 1: Análise de Potencial de Ação}
Um PA ventricular tem as seguintes características:
\begin{itemize}
    \item V$_{rest}$ = -85 mV
    \item V$_{peak}$ = +30 mV
    \item APD$_{90}$ = 300 ms
\end{itemize}
\begin{enumerate}
    \item Calcule a amplitude do PA
    \item Se a frequência cardíaca é 75 bpm, qual a porcentagem do ciclo em que a célula está refratária?
    \item Como uma droga que bloqueia I$_{Kr}$ afetaria o APD?
\end{enumerate}
\end{block>

\begin{block}{Questão 2: Concentrações Iônicas}
Calcule o potencial de equilíbrio para Ca$^{2+}$ a 37°C, dados:
\begin{itemize}
    \item [Ca$^{2+}$]$_o$ = 2 mM
    \item [Ca$^{2+}$]$_i$ = 0.1 $\mu$M
\end{itemize>
Use a equação de Nernst: $E = \frac{RT}{zF}\ln\frac{[X]_o}{[X]_i} = \frac{58}{z}\log_{10}\frac{[X]_o}{[X]_i}$ mV a 37°C
\end{block>
\end{frame>

\begin{frame}{Exercícios - Parte 2}
\begin{block}{Questão 3: Condutância de Canais}
Uma célula cardíaca tem g$_{Na}$ = 23 mS/cm$^2$. Durante o upstroke do PA:
\begin{itemize}
    \item V$_m$ = 0 mV
    \item E$_{Na}$ = +50 mV
    \item Ativação: m$^3$h = 0.9 $\times$ 0.6 = 0.54
\end{itemize>

Calcule a corrente I$_{Na}$ (em $\mu$A/cm$^2$).
\end{block>

\begin{block}{Questão 4: Acoplamento EC}
\begin{enumerate}
    \item Desenhe o fluxograma completo do acoplamento excitação-contração
    \item Explique como uma redução de 50\% em I$_{Ca,L}$ afetaria a contração
    \item Por que bloqueadores de canal de Ca$^{2+}$ são efetivos no tratamento de hipertensão?
\end{enumerate}
\end{block>

\begin{block}{Questão 5: Comprimento de Onda de Reentrada}
Se a velocidade de condução é $\theta$ = 0.5 m/s e o ERP = 250 ms:
\begin{enumerate}
    \item Calcule o comprimento de onda $\lambda$
    \item Qual o perímetro mínimo de circuito para sustentar reentrada?
\end{enumerate}
\end{block>
\end{frame>

\begin{frame}{Exercícios - Parte 3}
\begin{alertblock}{Questão 6: Projeto Computacional}
Implemente um modelo simplificado de cardiomiócito ventricular:
\begin{enumerate}
    \item Use o modelo fornecido nos slides (Noble ou LR simplificado)
    \item Simule um potencial de ação com estímulo em t = 5 ms
    \item Plote V$_m$(t), [Ca$^{2+}$]$_i$(t) e as principais correntes iônicas
    \item Investigue o efeito de:
    \begin{itemize}
        \item Redução de 50\% em g$_{Kr}$ (simulando bloqueio de hERG)
        \item Aumento de 20\% em g$_{K1}$ (simulando hiperkalemia)
        \item Redução de 30\% em g$_{Ca,L}$ (simulando bloqueador de canal)
    \end{itemize}
    \item Para cada manipulação, calcule APD$_{90}$ e compare com controle
\end{enumerate}
\end{alertblock>

\begin{block}{Questão 7: Análise Clínica}
Um paciente apresenta:
\begin{itemize}
    \item FC: 180 bpm (taquicardia)
    \item Ritmo irregular
    \item Ausência de ondas P
    \item Complexo QRS estreito
\end{itemize}
\begin{enumerate}
    \item Qual o diagnóstico mais provável?
    \item Explique o mecanismo fisiopatológico
    \item Quais opções terapêuticas você consideraria?
\end{enumerate}
\end{block}
\end{frame>

% ============================================================================
% SEÇÃO 8: REFERÊNCIAS
% ============================================================================
\section{Referências}

\begin{frame}{Referências Principais}
\begin{thebibliography}{99}
\bibitem{keener} Keener J., Sneyd J. (2009). \textit{Mathematical Physiology}, 2nd ed. Springer.

\bibitem{katz} Katz A.M. (2010). \textit{Physiology of the Heart}, 5th ed. Lippincott Williams \& Wilkins.

\bibitem{noble} Noble D. (1962). A modification of the Hodgkin-Huxley equations applicable to Purkinje fibre action and pacemaker potentials. \textit{J Physiol}.

\bibitem{luo} Luo C.H., Rudy Y. (1991). A model of the ventricular cardiac action potential. \textit{Circ Res}.

\bibitem{tentusscher} ten Tusscher K.H., Panfilov A.V. (2006). Alternans and spiral breakup in a human ventricular tissue model. \textit{Am J Physiol Heart Circ Physiol}.

\bibitem{oruru} O'Hara T., Virág L., Varró A., Rudy Y. (2011). Simulation of the undiseased human cardiac ventricular action potential. \textit{PLoS Comput Biol}.
\end{thebibliography}
\end{frame>

\begin{frame}{Leitura Complementar}
\begin{block}{Livros}
\begin{itemize}
    \item Clayton R.H. et al. (2011). \textit{Models of cardiac tissue electrophysiology}. Prog Biophys Mol Biol.
    \item Trayanova N.A. (2011). \textit{Whole-heart modeling}. Circ Res.
    \item Zipes D.P., Jalife J. (2013). \textit{Cardiac Electrophysiology: From Cell to Bedside}, 6th ed.
    \item Klabunde R. (2011). \textit{Cardiovascular Physiology Concepts}, 2nd ed.
\end{itemize}
\end{block>

\begin{block}{Recursos Online}
\begin{itemize}
    \item \textbf{CellML Repository:} \url{https://models.cellml.org/electrophysiology}
    \item \textbf{Cardiac Simulation:} \url{http://www.cardiac-sim.org}
    \item \textbf{OpenCARP:} \url{https://opencarp.org}
    \item \textbf{PhysioNet:} \url{https://physionet.org} (dados de ECG)
    \item \textbf{LivingHeart Project:} \url{https://www.3ds.com/simulia}
\end{itemize>
\end{block>
\end{frame>

% ============================================================================
% SLIDE FINAL
% ============================================================================
\begin{frame}
\begin{center}
{\Huge Obrigado!}

\vspace{1cm}

{\Large Capítulo 12: Physiological Systems - Heart}

\vspace{1cm}

\textbf{Próximo Capítulo:}\\
Physiological Systems - Neuronal Dynamics

\vspace{1cm}

\textit{Dúvidas? Perguntas?}
\end{center}
\end{frame}

\end{document}
